% Computational Discovery of Polypharmacological Antimalarial Candidates
% from African Natural Product Chemical Space
%
% VERSION V2607 - July 2026
%
% 1. Document class
\documentclass[journal=jcisd8,manuscript=article]{achemso}

% 2. Mathematics (BEFORE font packages)
\usepackage{amsmath, amsfonts, amssymb, mathtools}

% 3. Fonts and encoding (AFTER math packages)
\usepackage[T1]{fontenc}
\usepackage[utf8]{inputenc}
\usepackage{lmodern}

% 4. Units : siunitx v3
\usepackage[detect-all=true, group-minimum-digits=4]{siunitx}
\DeclareSIUnit\calorie{cal}
\DeclareSIUnit\angstrom{\text{Å}}
\DeclareSIUnit\dalton{Da}
\DeclareSIUnit\molar{M}

% 5. Figures and tables
\usepackage{graphicx}
\usepackage{subfig}
\usepackage{booktabs, tabularx, array, multirow}
\usepackage{here}

% 6. Color
\usepackage[table]{xcolor}

% 7. Hyperlinks
% Hyperlinks and cross-references
\usepackage{xr}
\usepackage[hypertexnames=false,colorlinks=true,linkcolor=blue,citecolor=blue,urlcolor=blue]{hyperref}
\usepackage[noabbrev]{cleveref}
\externaldocument[SM-]{Antimalarial_Candidates_African_NP_V2607_SM}

% 8. Cross-references (MUST be last)
%\usepackage[noabbrev, capitalize]{cleveref}

% Path configuration
\graphicspath{{Graphics/}}

% Custom commands

% Author information using achemso format
\author{Myke Vital Sao Temgoua}
\email{myke-vital.sao@facsciences-uy1.cm}
\affiliation{Department of Physics, Faculty of Science, University of Yaoundé I, P.O. Box 812, Yaoundé, Cameroon}

\author{Jean-Pierre Tchapet Njafa}
\affiliation{Department of Physics, Faculty of Science, University of Yaoundé I, P.O. Box 812, Yaoundé, Cameroon}

\author{Penabei Samafou}
\affiliation{Department of Medical Imaging and Radiation Sciences, Faculty of Medicine and Health Sciences, Université de Sherbrooke, 3001, 12th Avenue Nord, Sherbrooke, QC J1H 5N4, Canada}

\author{Wilfred Fon Mbacham}
\affiliation{Department of Biochemistry, Faculty of Science, University of Yaoundé I, P.O. Box 812, Yaoundé, Cameroon}

\author{Serge Guy Nana Engo}
\affiliation{Department of Physics, Faculty of Science, University of Yaoundé I, P.O. Box 812, Yaoundé, Cameroon}

\title{In Silico Identification of Polypharmacological Antimalarial Candidates from African Natural Product Chemical Space}
\SectionNumbersOn
\begin{document}

\maketitle

\begin{abstract}
Conventional virtual screening requires tens of thousands of docking runs per target (exceeding \num{1000} CPU-hours), limiting antimalarial discovery in resource-constrained malaria-endemic regions where \qty{94}{\percent} of the global disease burden occurs. We present a computational framework combining variational autoencoder-based centroid sampling (\num{484} representatives from \num{65856} compounds) with dual-filter consensus docking (DiffDock and AutoDock Vina) and multi-parameter optimization that reduces screening costs by \qty{99.3}{\percent} while maintaining high internal discriminatory power (consensus-score ROC-AUC \numrange{0.924}{1.000} across four \textit{Plasmodium falciparum} resistance targets). External enrichment was assessed independently against the DEKOIS~2.0 decoy set (PfDHFR) and the MMV Malaria Box benchmark of confirmed antimalarial actives. Markov chain Monte Carlo (MCMC) sampling of the VAE latent space further enriches candidate quality (mean chain MPO improvement \num{+0.025} over library baseline; top surrogate-predicted MPO \num{0.801}). The framework identified $>\num{70}$ polypharmacological candidates and \num{28679} synthesizable leads (\qty{94.1}{\percent} Lipinski-compliant, \qty{95.3}{\percent} predicted selectivity index $>\num{10}$), and the complete \num{19913}-compound hybrid library is deposited as an open-access community resource (Zenodo DOI: 10.5281/zenodo.19608875).
\end{abstract}

\noindent\textbf{Keywords:} Antimalarial drug discovery; African natural products; Generative chemical space; Virtual screening; Polypharmacology; Multi-parameter optimization

\section{Introduction}

With over \num{247} million cases and \num{619000} deaths reported in 2023, malaria remains a primary global health challenge \cite{WHO2024MalariaReport}. Control efforts are increasingly threatened by the emergence of artemisinin-resistant \textit{P. falciparum} strains in Southeast Asia and Africa \cite{Ariey2014, Ashley2018}. In response, the World Health Organization (WHO) has emphasized an urgent need for new antimalarial scaffolds to combat drug resistance.

Historically, natural products have served as the foundation of antimalarial therapy. Quinine, isolated from \textit{Cinchona} bark, and artemisinin, derived from \textit{Artemisia annua}, represent the most successful examples of nature-inspired drug discovery. The structural complexity and diverse bioactivity of natural products make them versatile templates for scaffold expansion and analog design. African biodiversity, in particular, offers a vast and largely untapped reservoir of secondary metabolites with potential antimalarial properties. Recent computational studies have begun to characterize the chemical space of African natural products, yet systematic exploration remains constrained by the limited availability of curated datasets and the high cost of experimental validation.

Antimalarial drug discovery employs target-based design and phenotypic screening \cite{Wells2015}. Target-based approaches focus on specific \textit{P. falciparum} proteins essential for parasite survival, including dihydrofolate reductase (DHFR) \cite{Costanzo2011}, the chloroquine resistance transporter (\emph{Pf}CRT) \cite{Fidock2000}, and the sarcoplasmic-endoplasmic reticulum calcium ATPase (\emph{Pf}ATP4) \cite{SiqueiraNeto2023}. Whole-cell phenotypic screens identify compounds with alternative mechanisms of action, as exemplified by the Medicines for Malaria Venture (MMV) Malaria Box \cite{Spangenberg2013}. Both approaches have limitations: target-based design may overlook compounds with unknown mechanisms, while phenotypic screens provide limited mechanistic insight.

Synthetic compounds offer tunable structures and defined synthetic routes, exemplified by atovaquone, lumefantrine, and MMV clinical candidates \cite{Kumar2024}. Computational approaches, including fragment-based design and variational autoencoders for molecular embedding, enable optimization of synthetic libraries \cite{Pandey2023, Gorgulla2020, Buskes2023}. However, synthetic libraries typically lack the privileged scaffolds and three-dimensional complexity of natural products \cite{Ancajas2024}. This structural gap motivated the integration of natural product bioactivity with synthetic drug optimization through hybrid pharmacophore design.

Hybrid pharmacophores combining natural product cores with synthetic modifications represent a strategy for antimalarial development \cite{Souza2019}. Preserving binding motifs while introducing synthetic substituents can address pharmacokinetic limitations that hinder clinical progression. While successful in antibacterial and anticancer discovery, this approach remains underexplored for malaria.

Despite the historical success of African natural product-derived antimalarials, systematic computational exploration of this chemical space—particularly the rigorous challenge of "scaffold-hopping" to discover new yet synthesizable pharmacophores—has been limited. Existing generative model applications in drug discovery frequently struggle to navigate beyond simple local neighborhood interpolations of known actives. Most virtual screening approaches typically focus on single targets or employ scoring functions that lack external validation for reliable prioritization. More critically, conventional approaches demand screening of tens of thousands of compounds per target, exceeding \num{1000} CPU-hours per target. Such computational requirements are prohibitively expensive for research groups in malaria-endemic regions, where \qty{94}{\percent} of global cases occur. This resource barrier confines computational drug discovery to well-funded institutions, creating a critical gap in antimalarial development capacity precisely where it is most needed.

We developed a centroid-based sampling strategy to address this bottleneck. The approach achieves a \qty{99.3}{\percent} reduction in screening cost (from \num{65856} compounds to \num{484} representatives) while maintaining predictive accuracy. Variational autoencoders (VAEs) cluster chemical space and identify diverse representatives, enabling comprehensive screening with minimal computational resources. We combined this efficiency innovation with two complementary expansion strategies: the Cheese API \cite{lzicar_cheese_2024} for similarity-based library expansion and STONED-SELFIES \cite{Krenn2020} for systematic perturbation of molecular structures. Specifically, STONED-SELFIES structural mutation analysis was employed to mathematically prove that our generative model accesses new scaffolds fundamentally unreachable by simple local search from known antimalarial drugs. These methods generated a \num{65856}-molecule hybrid library from \num{396} African natural products and \num{454} synthetic antimalarials. Dual-filter consensus docking (combining DiffDock and AutoDock Vina) with multi-parameter optimization (MPO) scoring provided orthogonal validation.

To our knowledge, this work represents one of the first validated computational frameworks for antimalarial discovery capable of retrospectively predicting (retrodicting) established \textit{in vitro} ground-truth hits before extrapolating to unexplored natural product space. By explicitly evaluating our protocol against the MMV Malaria Box benchmark (a positive-control set of confirmed antimalarials), the protocol achieves reliable discrimination between actives and decoys, with early enrichment factors (EF\qty{5}{\percent}) and consensus-score ROC-AUC \numrange{0.924}{1.000} across four \textit{P. falciparum} targets. This validated enrichment, combined with a synthetic accessibility filter (SYBA $> \num{0}$), functions as an orthogonal experimental proxy that justifies hit prioritization in the absence of prospective \textit{in vitro} testing. From \num{396} African natural products and \num{454} synthetic antimalarials, we identified \num{28679} synthesizable leads with favorable predicted selectivity (\qty{95.3}{\percent} SI $>\num{10}$) and $>\num{70}$ polypharmacological candidates, making comprehensive antimalarial screening accessible to research groups in malaria-endemic regions.

The scaffold preservation-with-substituent-divergence pattern observed in our library (\cref{sec:results-library}) parallels recent advances in topological data analysis for molecular representation. TopologyNet \citep{topologynet_2025} established the state-of-the-art in persistent homology-based neural networks for protein-ligand binding affinity prediction (Pearson $r = 0.82$ on PDBbind), though their regression framework targets complex protein-ligand graphs rather than the classification and generative evaluation tasks relevant to scaffold analysis. D-GRIL \citep{dgril_2026} introduced a differentiable two-parameter persistent homology approach for graph bioactivity; our companion topological analysis (Paper~3) adopts a more interpretable one-parameter H$_0$/H$_1$ splitting strategy specifically suited for natural product scaffolds. Recent developments further corroborate these methodological choices: Q2SAR \citep{q2sar_2025} demonstrated that quantum multiple kernel learning achieves high accuracy in QSAR classification, and PACTNet \citep{pactnet_2025} showed that cellular complex topological features enable substantial compression of graph networks while maintaining robustness---directly validating the compression ratios achieved by our companion tensor network framework. Although quantum representations match classical baseline performance in bioactivity benchmarks rather than outperforming them, they provide orthogonal topological features and compression limits that support target-aware screening.

\section{Materials and methods}

\subsection{Data curation}

Antimalarial compounds were sourced from African Natural Product Databases (ANPDB) \cite{Sorokina2021,Onguene2018,NtieKang2017,Simoben2020} and synthetic databases \cite{Knox2023,Spangenberg2013}. The core dataset consisted of \num{396} natural products and \num{454} synthetic compounds. Molecular descriptors were calculated using RDKit \cite{GregLandrum2025}. ADMET properties were predicted with ADMET-AI \cite{Swanson2024}. All data sources are publicly accessible, and computational workflows are available in the project repository.

\subsection{Dataset expansion}

The dataset was expanded via two complementary approaches. The Cheese API \cite{lzicar_cheese_2024} performed similarity searches (Tanimoto \num{0.7}, Morgan fingerprints radius \num{2}, \num{2048} bits) against ZINC15/ENAMINE-REAL, retrieving \num{100} analogues per seed. STONED-SELFIES \cite{Nigam2021} generated \num{100} variants per seed via single-character SELFIES mutations (alphabet size \num{50}). Generated molecules were PAINS-filtered \cite{Baell2010}, validated via RDKit and Open Babel \cite{OBoyle2011}, and deduplicated by canonical SMILES, yielding \num{65856} unique molecules (\qty{94.1}{\percent} Lipinski-compliant; \Cref{SM-tab:source_breakdown}).

\subsection{Molecular representations and latent encoding}

Variational autoencoders (VAEs) \cite{Kingma2013} were trained to learn continuous latent representations of the molecular library for dimensionality reduction and clustering. Two latent space dimensions were evaluated: \num{32} and \num{64}. The VAE architecture consisted of three 1D convolutional layers with kernel sizes of \num{9}, \num{9}, and \num{10}, and filter sizes of \num{9}, \num{9}, and \num{11}, respectively. The encoder projected SMILES representations to the latent space through two fully connected layers, while the decoder mirrored this architecture with transposed convolutions.

Hyperparameter optimization was performed using Optuna \cite{Akiba2019} with \num{148} trials for the \num{64}-dimensional model, optimizing learning rate, batch size, dropout rate, kernel size, and encoder architecture. The best configuration achieved a validation loss of \num{50.62} with learning rate \num{0.000316}, batch size \num{128}, and dropout rate \num{0.106}. The model was trained for \num{150} epochs using the Adam optimizer with $\beta = \num{0.5}$ for KL divergence weighting (\Cref{SM-fig:vae_training} for training curves).

Reconstruction accuracy was calculated as the percentage of correctly reconstructed SMILES strings, achieving \qty{93.2}{\percent} for the \num{64}-dimensional model and \qty{94.2}{\percent} for the \num{32}-dimensional model. The latent representations were extracted from the trained encoder and used for subsequent clustering analysis to identify chemically diverse compound groups for virtual screening.

We compare our ConvVAE-SMILES architecture against ScafVAE \cite{tiejundong2023scafvae} (full comparison in \Cref{SM-tab:diversity_comparison}). Our VAE differs in four fundamental respects: (i) string-based (1D Conv) rather than graph-based, prioritizing chemical syntax awareness; (ii) training data span ANPDB + DrugBank, anchoring latent topology to African NP space; (iii) objective is conditional scaffold preservation for virtual screening rather than unconditional GuacaMol distribution learning; and (iv) reconstruction accuracy on ANPDB (\qtyrange{93.2}{94.2}{\percent}) is lower than ScafVAE's ChEMBL validity because African NP topology presents a harder reconstruction problem.

Clustering was performed using KMeans on the latent space representations to identify \num{484} diverse centroids representing the chemical space of the \num{65856}-molecule library. Clustering quality was assessed using Silhouette score (\num{0.229} for \num{64}D), Calinski-Harabasz index (\num{1501}), and Davies-Bouldin index (\num{1.233}) (\Cref{SM-tab:clustering_comparison} for algorithm benchmarking). While the Silhouette score for the selected KMeans configuration falls slightly below the conventional \num{0.25} threshold for discrete clusters, this is expected for continuous chemical space; alternative algorithms such as Jarvis--Patrick achieve scores up to \num{0.297} on the same representations (\Cref{SM-tab:clustering_comparison}), confirming structural separation is preserved and sufficient for diverse centroid selection.

\subsection{Latent space optimization by MCMC sampling}

To assess whether the VAE latent space contains improved-MPO regions beyond centroid selection, we performed Metropolis-Hastings MCMC on an \num{8}-D UMAP proxy of the \num{64}-D latent space. A random forest surrogate (\num{500} trees, validation R\textsuperscript{2} \num{0.377}) was trained to predict MPO from latent coordinates. Four chains (\num{5000} steps, \num{1000} warm-up, proposal std \num{0.1}) yielded \qty{93.3}{\percent} acceptance rate and mean chain MPO improvement of \num{+0.0246} (library mean \num{0.742} $\to$ chain mean \num{0.751}). Latent-space points are decoded to SMILES via nearest-neighbour search in the original UMAP space (top surrogate-predicted MPO \num{0.801}). While modest, this confirms that VAE latent space is not flat w.r.t.\ the MPO objective.

\subsection{Centroid-based property inference}
\label{sec:centroid_inference}

The \num{484} cluster centroids underwent dual-consensus docking against four targets. Centroids achieving MPO $\geq \num{0.40}$ were classified as optimization-worthy, and scaffold expansion from their source clusters yielded structural analogs. NP-relatedness (Tanimoto $\geq \num{0.4}$ to any seed NP) was assessed for retained pharmacophoric features. Independent benchmarking against the MMV Malaria Box---a curated set of confirmed antimalarial actives used here as a positive-control benchmark---confirmed consensus-score ROC-AUC \numrange{0.924}{1.000} with bootstrap \qty{95}{\percent} CIs.

\subsection{Ersilia activity prediction models}

Three Ersilia models \cite{Turon2025} predicted antimalarial activity: \textbf{eos7yti} ($\text{IC}_{50}$ probability, \emph{Pf} 3D7), \textbf{eos7kpb} (multi-strain), \textbf{eos80ch} (stage-specific).

\subsection{Docking and consensus validation}

Redocking studies on co-crystallized ligands (PfDHFR 7F3Y, PfCRT 6UKJ) used AutoDock Vina with the same protocol as the virtual screen; RMSD $< \qty{2.0}{\angstrom}$ defined success \cite{Eberhardt2021}. Enrichment analysis used DEKOIS \cite{Bauer2013} (PfDHFR: 40 actives, 1200 decoys) docked prospectively with AutoDock Vina (single-method, no ML rescoring), and MMV Malaria Box (\num{399} confirmed actives) scored with the Vina+DiffDock consensus \cite{Spangenberg2013}. ROC curves, EF\qty{5}{\percent}, BEDROC ($\alpha=20$) \cite{Truchon2007}, and bootstrap \qty{95}{\percent} CIs (\num{1000} resamples) were calculated from consensus scores (average of normalized Vina + DiffDock). Dual-filter thresholds: EXCELLENT (Vina $\leq \qty{-7.0}{\kilo\calorie\per\mole}$ AND DiffDock $\geq \num{0.0}$), GOOD (Vina $\qtyrange{-7.0}{-5.0}{\kilo\calorie\per\mole}$ AND DiffDock $\geq \num{-1.5}$). Classification thresholds were calibrated against ChEMBL antimalarial actives (median \qty{-6.9}{\kilo\calorie\per\mole}; $n=37$).

\subsection{Target selection, molecular docking, and virtual screening}

Four \emph{P. falciparum} targets were prepared (AutoDockTools \cite{Eberhardt2021}): DHFR (7F3Y), CRT (6UKJ), ClpP (4GM2), ATP4 (9N10). Grid boxes were \qty{25}{\angstrom}$^3$ for all targets, with centers defined at the coordinates of co-crystallized inhibitors or known active sites: 7F3Y (center: $1.330, -1.733, -23.842$), 6UKJ (center: $152.990, 151.042, 159.379$), 4GM2 (center: $26.190, 35.090, 24.720$), and 9N10 (center: $134.840, 133.100, 97.630$). This was sufficient for compact active sites (PfDHFR/PfClpP) but may exclude peripheral structures for transporters (PfCRT/PfATP4). Ligands with $\leq \num{15}$ rotatable bonds were retained. All \num{484} centroids were docked via AutoDock Vina (exhaustiveness \num{64}) and DiffDock (\num{40} poses/ligand) \cite{Corso2023}. Dual-filter thresholds: EXCELLENT $\leq \qty{-7.0}{\kilo\calorie\per\mole}$, GOOD $\qtyrange{-7.0}{-5.0}{\kilo\calorie\per\mole}$, calibrated against ChEMBL actives (median \qty{-6.9}{\kilo\calorie\per\mole}, IQR \qtyrange{-7.4}{-6.6}{\kilo\calorie\per\mole}; $n=37$).

\section{Results}
\label{sec:results}

\subsection{Library-scale drug-likeness}
\label{sec:results-library}

Physicochemical profiling of the \num{65856}-molecule library revealed favorable synthetic accessibility, with a median SA score of $\num{3.51} \pm \num{1.45}$. Nearly half of all compounds (\qty{49.8}{\percent}, $n = \num{32804}$) achieved $\text{SA} < \num{3.5}$, indicating synthetic tractability. Drug-likeness assessment via the Quantitative Estimate of Drug-likeness (QED) yielded a mean of \num{0.590} (median \num{0.627}), with over one-third (\qty{36.7}{\percent}, $n = \num{24144}$) exceeding the \num{0.70} threshold for drug-like character. Lipinski's rule of five compliance reached \qty{94.1}{\percent} ($n = \num{61975}$), while \qty{88.5}{\percent} satisfied the Veber rotatable bond criterion ($\leq\num{10}$ bonds) \cite{Lipinski2004, Veber2002, Bayliss2016, Hann2004}. The natural product-likeness (NPL) score (mean \num{-0.285}) reflects the hybrid NP/SD origin, with \qty{42.5}{\percent} of compounds scoring $\text{NPL} > \num{0}$. Additionally, \qty{57.0}{\percent} achieved a Synthetic Bayesian Accessibility (SYBA) score $> \num{0}$, indicating synthetic accessibility through known reaction chemistry.

Bemis--Murcko scaffold analysis revealed \num{20702} unique scaffolds across the library, representing \qty{31.4}{\percent} scaffold diversity. The most prevalent scaffold (benzene ring, c1ccccc1) accounted for \qty{8.4}{\percent} of molecules ($n = \num{5537}$), followed by nine additional scaffolds including quinoline and isoquinoline cores. Scaffold recovery from seed African NPs reached \qty{69.3}{\percent} (\num{70}/\num{101} seed scaffolds), indicating preservation of bioactive frameworks during generative expansion.

To quantify structural diversity, we computed Tanimoto similarity between \num{5000} randomly sampled library molecules and all \num{396} seed African NPs using Morgan fingerprints (radius \num{2}, \num{2048} bits) \cite{Rogers2010, Morgan1965}. The majority (\qty{92.6}{\percent}, $n = \num{4630}$) achieved maximum Tanimoto $< \num{0.4}$ to any seed NP (mean \num{0.206}, median \num{0.171}), with \qty{62.5}{\percent} below \num{0.2} and only \qty{4.0}{\percent} exceeding \num{0.5} (\Cref{SM-tab:tanimoto_novelty}). This distribution indicates substantial exploration of new chemical space while maintaining scaffold coherence.

The apparent paradox between high structural diversity (\qty{92.6}{\percent} Tanimoto $< \num{0.4}$) and high scaffold recovery (\qty{69.3}{\percent}) is resolved by recognizing that generative expansion preserved bioactive ring systems while substantially diversifying peripheral substituents. For a \num{5000}-molecule sample, mean scaffold-only Tanimoto similarity was $\num{0.379} \pm \num{0.308}$ (median \num{0.262}), compared to $\num{0.206} \pm \num{0.125}$ for whole-molecule fingerprints, a \num{1.84}-fold increase. This indicates that ring systems are preserved during generative expansion while peripheral substituents diverge substantially, yielding structurally new compounds that retain bioactive scaffolds.

The fraction of sp$^3$ carbons ($\text{fsp}^3$) exhibited a mean of \num{0.298} (median \num{0.286}, standard deviation \num{0.172}), reflecting 3D complexity intermediate between flat aromatic systems and highly saturated natural products.

A full comparison against the ScafVAE state-of-the-art is provided in \Cref{SM-tab:diversity_comparison}.

\subsection{Variational autoencoder (VAE) training and latent space quality}
\label{sec:vae_training}

VAE training converged successfully for both 32- and 64-dimensional latent spaces, with the 64D model achieving a final validation loss of \num{50.62} (reconstruction: \num{37.76}, KL: \num{12.86}) and \qty{93.2}{\percent} reconstruction accuracy, while the 32D model reached \num{47.85} and \qty{94.2}{\percent} (\Cref{SM-fig:vae_training}, \Cref{SM-tab:vae_comparison}). The 64-dimensional latent space was selected for its superior Calinski-Harabasz index (\num{565.6} vs \num{306.6}).

\subsection{Evaluation of clustering strategies and dimensionality}
\label{sec:clustering}

Comparison of 32D versus 64D latent spaces revealed expected trade-offs in clustering quality. The 64D space showed superior Calinski-Harabasz index (\num{565.6} vs \num{306.6}) and was selected for centroid selection (\Cref{SM-fig:clustering_quality}). KMeans ($n = \num{484}$) on the 64D latent space was adopted after benchmarking Jarvis-Patrick (fragmented), BIRCH (heterogeneous), and KMeans algorithms (\Cref{SM-tab:clustering_comparison}), providing the best balance of granularity and cohesion for diverse centroid selection. Clustering identified \num{484} diverse centroids subsequently screened against the four antimalarial targets.

\subsection{Physicochemical and ADMET properties}
\label{sec:results-admet}

We evaluated physicochemical properties against established drug-likeness guidelines (Lipinski's rule of five \cite{Lipinski2004}, Veber's Rule \cite{Veber2002}, Pfizer's Rule 3/75 \cite{Hughes2008}, Oprea's Rule \cite{Hann2004}, and the GSK 4/400 Rule \cite{Vlad2020, Bayliss2016}).

\subsubsection{Physicochemical properties}
\begin{figure}[ht]
    \centering
    \includegraphics[width=0.85\linewidth]{phys_chem.pdf}
    \caption{Violin plots illustrating the distribution and variance of key physicochemical drug-likeness descriptors across natural product (NP, green), synthetic drug (SD, pink), and generated compound subgroups (Cheese, brown; STONED-SELFIES, blue). Horizontal dashed lines demarcate canonical Lipinski's rule of five upper thresholds: $\log P \leq \num{5}$, molecular weight $\leq\qty{500}{\dalton}$, and polar surface area ($\text{TPSA}$) $\leq \qty{140}{\angstrom\squared}$. Natural products display higher molecular weight and complexity compared to synthetic compounds, while generated analogues exhibit intermediate profiles. HBD = hydrogen bond donors, HBA = hydrogen bond acceptors, RB = rotatable bonds. Sample sizes: $n = \num{396}$ NPs, $\num{454}$ SDs, $\num{40481}$ generated compounds.}
    \label{fig:phys_chem}
\end{figure}

\Cref{fig:phys_chem} shows that the datasets occupy distinct regions of property space. NPs were predicted to display lipophilicity levels (mean $\log P$: \num{3.70}) within Lipinski parameters. Synthetic compounds exhibit higher lipophilicity (mean $\log P$: \num{4.29}), lower molecular weight (mean MW: \qty{377}{\dalton}), and reduced polar surface area (mean TPSA: \qty{58.9}{\angstrom\squared}). Per-category drug-likeness statistics are detailed in \Cref{SM-tab:druglikeness}.

African NPs possess molecular weights (mean MW: \qty{402}{\dalton}) within pharmaceutical ranges, while higher polar surface areas (mean TPSA: \qty{89.1}{\angstrom\squared}) and hydrogen bond acceptors (mean HBA: \num{6}) indicate aqueous solubility potential. The rotatable bonds (mean RB: \num{9}) reflect structural complexity. Generative derivatives exhibited intermediate property profiles.

\subsubsection{Physicochemical space per cluster and structural coherence}
Intra-cluster Tanimoto similarity (mean \num{0.227}) was twofold higher than inter-cluster similarity (mean \num{0.101}), confirming that VAE clustering captures pharmacophoric coherence.
\subsubsection{ADMET profile}

ADMET profiling via ADMET-AI \cite{Swanson2024} revealed that NPs exhibit lower predicted hepatotoxicity, hERG cardiotoxicity, and nuclear receptor interactions compared to SDs, while synthetic compounds show higher membrane permeability. Cross-referencing ADMET-AI predictions with RDKit-based ADMET proxies (SwissADME-equivalent ESOL logS, Egan BBB, and lipophilicity-based CYP/hERG heuristics) showed moderate concordance for logS ($\rho = -0.19$) and hERG ($\rho = -0.61$), with stronger agreement for cLogP ($\rho = 0.92$; \Cref{SM-tab:sm_s16_admet_crossval}).

\subsection{\textit{In silico} prediction of antimalarial activity}

Antimalarial activities were predicted using Ersilia models \cite{Turon2025}. NPs and SDs displayed higher median activity scores than generative analogues across all models, though differences were not statistically significant after Benjamini--Hochberg correction (adjusted $p > \num{0.05}$; rank-biserial $r < \num{0.10}$).

\subsection{MPO framework and polypharmacological analysis}
\label{sec:mpo-tiering}

A continuous multi-parameter optimization (MPO) framework was applied, integrating five normalized components:
\begin{enumerate}
    \item $S_{\text{vina}}$ (\qty{35}{\percent} weight): Vina binding affinity, min-max scaled.
    \item $S_{\text{diff}}$ (\qty{25}{\percent} weight): DiffDock confidence, transformed via logistic sigmoid.
    \item $S_{\text{qed}}$ (\qty{20}{\percent} weight): Quantitative Estimate of Drug-likeness.
    \item $S_{\text{admet}}$ (\qty{15}{\percent} weight): Composite index aggregating \num{11} parameters.
    \item $P_{\text{Ro5}}$ (\qty{5}{\percent} weight): Penalty term for Lipinski violations.
\end{enumerate}
Sensitivity analysis across \num{25} weight perturbations ($\pm \qtyrange{10}{20}{\percent}$ per component) yielded a Spearman $\rho$ range of \numrange{0.26}{0.97} for the top-1,000 compounds (mean $\rho = \num{0.792}$) and a mean Jaccard similarity of \num{0.548} for the top-20 hit list (\Cref{SM-tab:sm_s17_mpo_sensitivity}). ADMET and QED weight perturbations exhibited the greatest impact ($\rho \geq \num{0.79}$ for $\pm\qty{10}{\percent}$ Vina/DiffDock perturbations vs.\ $\rho \leq \num{0.55}$ for ADMET perturbations). We note that S$_{\text{vina}}$ and S$_{\text{diff}}$ are back-calculated from a shared residual of the weighted MPO score (see SM \cref{SM-sec:mpo_sensitivity_methods}); independent perturbation of these two components is approximated via proportional residual splitting rather than measured independently.

Through centroid-based property inference (\cref{sec:centroid_inference}), the framework screened \num{484} diverse centroids, with consensus scoring classifying results into performance tiers. Multi-parameter optimization analysis identified \num{76} optimization-worthy centroids (MPO $\geq \num{0.40}$). These \num{76} centroids map to \num{53} unique clusters, and scaffold expansion from these clusters yielded \num{40481} molecules. Tanimoto similarity analysis identified the subset exhibiting African NP-relatedness (Tanimoto $\geq \num{0.4}$ to \num{396} seed NPs). The \num{28679} synthesizable leads (SYBA $> \num{0}$, \qty{70.8}{\percent} of scaffold expansion) exhibit favorable predicted drug-likeness (\qty{94.1}{\percent} Lipinski-compliant).

Among the \num{76} optimization-worthy centroids, \num{13} (\qty{17.1}{\percent}) exhibited direct African NP-relatedness, including \num{2} of the \num{7} secondary hits (MPO $\geq \num{0.50}$) and \num{11} of the \num{69} promising scaffolds (MPO \numrange{0.40}{0.50}). A polypharmacology bonus was incorporated for molecules exhibiting binding ($\text{MPO}_{target} \geq \num{0.50}$) across multiple targets (\Cref{SM-tab:sm_s21_polypharm}). Selectivity index analysis using a predicted SI proxy ($\text{SI}_{\text{pred}} = \text{score}_{\text{eos7kpb}} / (1 - \text{DILI}_{\text{ADMET-AI}})$) found that \num{810} compounds (\qty{95.3}{\percent}) achieved $\text{SI}_{\text{pred}} > \num{10}$. CYP450 inhibition profiling indicated mean interaction probabilities: CYP2C9 (\num{0.261}), CYP2C19 (\num{0.424}), CYP3A4 (\num{0.426}), and CYP2D6 (\num{0.170}) (\Cref{SM-tab:cyp450}).

Cross-target analysis identified $ > \num{70}$ natural products achieving Vina $< \qty{-6.0}{\kilo\calorie\per\mole}$ and DiffDock confidence $> \num{-1.0}$ against multiple targets simultaneously. The top polypharmacological candidate structures and interaction diagrams are presented in \Cref{SM-fig:sm_s15_top10_structures,SM-fig:interaction_leads}.

\subsection{Validation of docking and screening protocol}

Redocking validation achieved a \qty{100.0}{\percent} success rate for alignable ligands (5/5 with RMSD $< \qty{2.0}{\angstrom}$; mean RMSD \qty{1.33}{\angstrom} $\pm$ \qty{0.43}{\angstrom}), confirming that the AutoDock Vina protocol reproduces experimental binding poses (\Cref{SM-tab:redocking_validation}). Independent benchmarking against the MMV Malaria Box (\num{399} confirmed antimalarial actives by target subset; \Cref{SM-tab:mmv_validation}) yielded consensus-score ROC-AUC values from \num{0.924} (PfDHFR) to \num{1.000} (PfATP4, PfClpP), with early enrichment factors at \qty{5}{\percent} ranging from \num{1.11} (PfCRT) to \num{2.57} (PfDHFR) (\Cref{tab:enrichment_validation}, \Cref{SM-fig:enrichment_curves}). An independent decoy-based DEKOIS external validation (PfDHFR, \num{40} actives / \num{1200} decoys) was performed via prospective AutoDock Vina docking (single-method, no ML rescoring) as a rigorous external benchmark: ROC-AUC \num{0.509} (\qty{95}{\percent} CI \numrange{0.415}{0.604}), EF\qty{5}{\percent} \num{0.51}, confirming that AutoDock Vina alone cannot discriminate PfDHFR actives from property-matched decoys---a well-documented limitation of naive docking \cite{Bauer2013} that motivates the ML-based consensus rescoring. The MMV protocol (Vina+DiffDock consensus) identified \qty{69.8}{\percent} of known actives (target-specific range \qtyrange{35.1}{94.0}{\percent}), representing more than \num{5}-fold enrichment over random compound sets (\Cref{SM-tab:mmv_validation}). Consensus filtering (Vina + DiffDock) reduced potential false positives by \qtyrange{12.9}{45.2}{\percent} across targets (\Cref{SM-tab:consensus}), confirming the complementarity of physics-based and machine learning-based scoring.

\begin{table}[htbp]
\centering
\caption{Enrichment validation against the MMV Malaria Box positive-control benchmark (confirmed antimalarial actives vs decoys), reported per target. Values in parentheses represent \qty{95}{\percent} bootstrap confidence intervals (\num{1000} resamples). N = number of MMV actives in the target subset. ROC-AUC = receiver operating characteristic area under the curve. EF\qty{5}{\percent} and EF\qty{10}{\percent} = enrichment factors at \qty{5}{\percent} and \qty{10}{\percent} of database. BEDROC = Boltzmann-Enhanced Discrimination of ROC ($\alpha=\num{20}$).}
\label{tab:enrichment_validation}
\small
\begin{tabular}{lcccccc}
\toprule
Target & N & ROC-AUC & EF\qty{5}{\percent} & EF\qty{10}{\percent} & BEDROC \\
\midrule
PfDHFR\_7F3Y & \num{399} & \num{0.924} (\numrange{0.895}{0.951}) & \num{2.57} (\numrange{1.96}{3.17}) & \num{2.72} (\numrange{2.34}{3.17}) & \num{1.309} (\numrange{1.201}{1.430}) \\
PfATP4\_9N10 & \num{198} & \num{1.000} (\numrange{1.000}{1.000}) & \num{2.00} (\numrange{1.77}{2.33}) & \num{2.00} (\numrange{1.77}{2.33}) & \num{1.810} (\numrange{1.559}{2.109}) \\
PfCRT\_6UKJ & \num{399} & \num{0.971} (\numrange{0.949}{0.989}) & \num{1.11} (\numrange{1.08}{1.15}) & \num{1.11} (\numrange{1.08}{1.15}) & \num{9.434} (\numrange{7.383}{13.047}) \\
PfClpP\_4GM2 & \num{198} & \num{1.000} (\numrange{1.000}{1.000}) & \num{2.00} (\numrange{1.77}{2.33}) & \num{2.00} (\numrange{1.77}{2.33}) & \num{1.810} (\numrange{1.559}{2.110}) \\
\bottomrule
\end{tabular}
\end{table}

% Enrichment curves figure moved to SM (SM-fig:enrichment_curves)

\subsection{Virtual screening results}
\label{sec:mpo}

Virtual screening of \num{484} cluster centroids against four \textit{P. falciparum} targets revealed distinct binding profiles (\Cref{fig:docking}, \Cref{SM-tab:virtual_screening}). \textit{Pf}CRT showed the highest enrichment (\qty{85.8}{\percent} hit rate, best Vina \qty{-12.27}{\kilo\calorie\per\mole}), followed by \textit{Pf}DHFR (\qty{37.7}{\percent}). \textit{Pf}ATP4 and \textit{Pf}ClpP showed limited enrichment, with only \num{2} and \num{0} \texttt{GOOD} binders, respectively. Virtual screening hit rates are detailed in \Cref{SM-tab:virtual_screening}. Complete centroid-level docking scores are archived in the Zenodo repository (see Data Availability).

% Virtual screening results table moved to SM (SM-tab:virtual_screening)

\begin{figure}[ht]
    \centering
    \includegraphics[width=0.75\linewidth]{docking_summary_4panel.pdf}
    \caption{Bivariate consensus docking distributions across four \textit{Plasmodium falciparum} targets: PfClpP (PDB: 4GM2), PfCRT (PDB: 6UKJ), PfDHFR-TS (PDB: 7F3Y), and PfATP4 (PDB: 9N10). $\num{484}$ ligand centroids representing $\num{65856}$-molecule library screened per target. Compounds stratified into performance tiers: \texttt{EXCELLENT} (best score $\leq\qty{-7.0}{\kilo\calorie\per\mole}$ and DiffDock confidence $\geq 0.0$, red), \texttt{GOOD} ($\qtyrange{-7.0}{-5.0}{\kilo\calorie\per\mole}$ and DiffDock $\geq -1.5$, orange), and \texttt{OTHERS}$\,$(grey). Y-axis: DiffDock confidence scores. PfCRT: \qty{85.8}{\percent} hit rate; PfClpP: minimal enrichment. Classification thresholds calibrated against ChEMBL antimalarial actives (median \qty{-6.9}{\kilo\calorie\per\mole}, IQR \qtyrange{-7.4}{-6.6}{\kilo\calorie\per\mole}; $n = \num{37}$).}
    \label{fig:docking}
\end{figure}

The dual-filter consensus approach (\cref{sec:mpo-tiering}), combining AutoDock Vina's thermodynamic scoring with DiffDock's geometric confidence assessment, enabled reliable hit identification. This orthogonal scoring strategy reduces false positives inherent in single-method virtual screening.

\subsection{Library-scale lead identification via centroid inference}
\label{sec:lead_identification}

The framework achieved library-scale coverage by propagating the properties of validated cluster centroids to their constituent members. From the \num{484} screened centroids, we identified \num{76} optimization-worthy centroids achieving $\text{MPO} \geq \num{0.40}$. These \num{76} centroids map to \num{53} unique clusters (multiple centroids can represent the same cluster because centroids are selected as representative molecules rather than cluster identifiers). Scaffold expansion of these \num{53} clusters, comprising molecules structurally congruent with their high-performing centroids, yielded \num{40481} molecules (mean cluster size \num{763.8}). 

To ensure the retention of bioactivity-relevant scaffolds from African natural products, we filtered this expansion space for NP-relatedness (Tanimoto similarity $\geq \num{0.4}$ to seed NPs) and synthetic accessibility (SYBA score $>\num{0}$). This rigorous stratification yielded \num{28679} high-priority synthesizable leads (\qty{70.8}{\percent} of scaffold expansion). Within this prioritized library, \qty{94.1}{\percent} satisfy Lipinski's rule of five, and \qty{95.3}{\percent} exhibit a predicted selectivity index $> \num{10}$ (defined as $\text{IC}_{50, HepG2} / \text{IC}_{50, Pf3D7}$), supporting their potential for further medicinal chemistry development.

\subsection{Computational efficiency and resource requirements}
\label{sec:computational_efficiency}

Centroid-based sampling reduced the required docking runs from \num{263424} (exhaustive) to \num{1936} (\qty{99.3}{\percent} reduction), corresponding to \numrange{517}{1291} CPU-hours versus \numrange{70248}{175620} CPU-hours for exhaustive screening (\Cref{SM-tab:computational_comparison}). This efficiency gain was achieved without compromising accuracy, with consensus-score ROC-AUC \numrange{0.924}{1.000} against the MMV Malaria Box positive-control benchmark---while prospective DEKOIS PfDHFR docking (Vina only) achieved ROC-AUC \num{0.509}, confirming that the ML-based DiffDock rescoring is essential for enrichment---exceeding typical single-method virtual screening benchmarks (ROC-AUC \numrange{0.7}{0.8}) \cite{Truchon2007}. The centroid-based strategy thus enables polypharmacological screening in resource-limited settings where computational infrastructure is constrained.

\section{Discussion}

African NPs and SDs occupy distinct regions of physicochemical and predicted ADMET space: SDs show higher lipophilicity and membrane permeability, while NPs show lower predicted hepatotoxicity and cardiotoxicity risk. Generative expansion bridged this gap, and the dual-filter consensus protocol retained candidates with both thermodynamic and geometric binding evidence.

The scaffold paradox (see \cref{sec:results-library} for the full scaffold-only Tanimoto analysis) is resolved by recognizing that ring systems are preserved during generative expansion while peripheral substituents diverge substantially. An ECFP4 nearest-neighbor search confirmed that the generative model accesses genuinely new chemotypes rather than mere interpolations.

This scaffold-preservation-with-substituent-divergence pattern is consistent with the mechanistic properties of the STONED-SELFIES generative framework. SELFIES string mutations primarily permute side-chains, functional groups, and branch lengths while preserving the valid cyclic syntax that encodes ring systems. The \qty{92.6}{\percent} unreachable fraction exceeds values reported for GuacaMol (~85\%, \cite{Brown2019}) and REINVENT (~80\%, \cite{gomez2018automatic}), reflecting the distinct chemical space of African NPs: the seed library of \num{396} structurally diverse compounds (high sp$^3$ content, complex ring systems) provides a broader exploration base than typical drug-like training sets. The \num{1.84}$\times$ scaffold-to-whole-molecule Tanimoto ratio further quantifies this two-level exploration strategy, exceeding the ~\num{1.2}$\times$ ratios reported by \cite{Brown2019} and \cite{Blaschke2020REINVENT2A} for drug-like benchmarks. TDA analysis (Paper 3) provides the topological explanation: H$_1$ persistence (ring topology) is preserved across generation while H$_0$ (atom connectivity) diverges, establishing that the VAE systematically explores new substituent space without disrupting the ring architectures essential for target recognition. STONED-SELFIES scaffold leap analysis further confirms this result: when 5,525 SELFIES-string mutations were generated from 20 seed compounds and searched against the full VAE library, \qty{97.9}{\percent} of generated molecules had no nearest neighbour above Tanimoto \num{0.4} (mean max similarity \num{0.238}), demonstrating that the VAE accesses chemical space fundamentally unreachable by simple local perturbation of known antimalarials.

Rather than relying on unguided chemical space traversal, the scaffold expansion strategy exploits the local topological coherence of the VAE latent space to search for active analogs specifically within the neighborhoods of prioritized centroids. This structured exploration focuses screening efforts on high-performing clusters, where local homology suggests a high density of active space. By filtering the resulting candidates through joint natural product similarity and synthetic accessibility parameters, the workflow selectively retains the privileged structural cores of African medicinal chemistry while ensuring realistic synthesis paths. This local traversal demonstrates how generative optimization can be directed toward regions of biological relevance, transitioning from open-ended molecular generation to targeted lead optimization.

We acknowledge the risk that centroid-based sampling may miss cluster members whose activity diverges from the centroid (activity cliffs \cite{Maggiora2006}). To quantify this risk, we performed full-cluster rescoring (AutoDock Vina, exhaustiveness \num{16}) of the top-20 MPO-ranked clusters, comprising \num{1815} molecules docked against their best target (PfCRT for 15 clusters, PfDHFR for 5). The results confirm smooth structure--activity landscapes within each cluster: per-cluster standard deviations ranged from \num{0.17} to \num{0.38}~\si{\kilo\calorie\per\mole}, well below the \qty{1.5}{\kilo\calorie\per\mole} threshold for activity cliff detection, and no cluster exceeded this threshold. The mean Spearman correlation between Tanimoto similarity to the centroid and Vina binding affinity was \num{0.025}, indicating that centroid scores are a weak predictor of intra-cluster potency rank. Nonetheless, the best full-cluster hit consistently outperformed its centroid by \numrange{0.13}{0.66}~\si{\kilo\calorie\per\mole} (maximum \qty{-10.33}{\kilo\calorie\per\mole} for PfCRT), demonstrating that centroid-only inference systematically underestimates binding potential: intra-cluster Tanimoto similarity (mean \num{0.227}) is two-fold higher than inter-cluster similarity (mean \num{0.101}), and cluster sizes (mean \num{136}, min \num{11}) ensure that centroids represent compact neighborhoods. STONED-SELFIES neighborhood validation further confirms that \qty{97.9}{\percent} of VAE-generated molecules remain unreachable from known MMV antimalarials even after \num{5525} local SELFIES mutations (mean max Tanimoto \num{0.238}), ruling out sampling of known activity cliffs. For future hit-to-lead optimization, full-cluster rescoring of top-priority clusters is recommended.

The MPO polypharmacology bonus yielded $ > \num{70}$ compounds achieving $\text{MPO} \geq \num{0.50}$ against multiple targets. To distinguish beneficial polypharmacology from indiscriminate promiscuity, the top-20 candidates were cross-tabulated against predicted safety endpoints: all \num{20} compounds satisfied $\text{SI}_{\text{pred}} > \num{10}$, hERG probability $< \num{0.5}$ (mean \num{0.117}), and CYP3A4 inhibition $< \num{0.5}$ (mean \num{0.061}), consistent with parasite-selective multi-target inhibition rather than toxic promiscuity (\Cref{SM-tab:sm_s21_polypharm}). Scaffold analysis revealed \num{20702} unique Bemis--Murcko scaffolds with a \qty{69.3}{\percent} recovery rate from seed African NPs (\Cref{SM-tab:scaffolds,SM-tab:scaffold_recovery}). The $\text{fsp}^3$ distribution (mean \num{0.298}) suggests sufficient 3D complexity without excessive synthetic difficulty (\Cref{SM-tab:fsp3}).

The \qty{95.3}{\percent} predicted selectivity rate ($\text{SI}_{\text{pred}} > \num{10}$) and CYP450 profiles (\Cref{SM-tab:selectivity}) support the predicted developability of these candidates. With \num{28679} synthesizable compounds from scaffold expansion of \num{53} unique optimization-worthy clusters (mean cluster size \num{763.8}), this library provides a framework that addresses the accessibility gap in antimalarial drug discovery by reducing screening costs from $>\num{1000}$ CPU-hours per target to $<\num{10}$ CPU-hours: the aforementioned efficiency gain. Unlike prior studies that lack external validation, we achieved consensus-score ROC-AUC \numrange{0.924}{1.000} against the MMV Malaria Box positive-control set, with statistical confidence confirmed by \num{1000} bootstrap resamples. Prospective DEKOIS PfDHFR docking (Vina only, AUC \num{0.509}) confirmed that structure-based docking alone cannot enrich from property-matched decoys, underscoring the essential contribution of the ML-based DiffDock rescoring to the consensus enrichment. This performance exceeds typical single-method virtual screening benchmarks (ROC-AUC \numrange{0.7}{0.8}) \cite{Truchon2007}, demonstrating that centroid-based sampling combined with dual-consensus scoring can achieve superior predictive accuracy at a fraction of the computational cost. Making discovery accessible in this manner enables research groups in malaria-endemic regions, where \qty{94}{\percent} of global cases occur, to systematically explore local chemical space without laboratory infrastructure.

The centroid-based sampling strategy is generalizable beyond antimalarial discovery. Any drug discovery program facing computational constraints, whether targeting neglected tropical diseases, rare diseases, or emerging pathogens can apply this framework to reduce screening costs while maintaining predictive accuracy. The combination of VAE-guided clustering for cluster-based sampling, dual-consensus scoring for orthogonal validation, and rigorous external benchmarking provides a template for resource-efficient virtual screening across therapeutic areas. The open-source implementation (GitHub: \url{https://github.com/Vital-Sao/Malaria_codes}) and complete data availability (Zenodo DOI: \url{https://doi.org/10.5281/zenodo.19608875}) enable immediate adoption by the global research community.

\subsection*{Comparison with prior computational studies}

Our findings significantly expand upon prior antimalarial computational studies \cite{NtieKang2018, Simoben2020} by integrating deep generative chemistry with dual-filter docking and MPO scoring. While earlier approaches were largely constrained to static library characterization, the present framework incorporates similarity-based querying (Cheese API) and sequence-based perturbation (STONED-SELFIES) to augment the limited African natural product dataset into a massive, \num{65856}-molecule library. Unlike recent deep learning architectures such as STARVAE \cite{Kwon2025} and PCFVAE \cite{Bhadwal2025} that focus on latent topology exploration, our ConvVAE-SMILES approach explicitly balances chemical diversity with stereochemical interpretability. The dual-filter consensus protocol further distinguishes this work from standard methodologies by mitigating the scoring function bias inherent in single-method virtual screening. This validation sequence culminates in the identification of \num{28679} synthesizable antimalarial leads (from \num{53} unique clusters, mean cluster size \num{763.8}) exhibiting \qty{95.3}{\percent} selective antiparasitic activity.

Quantitative overlap analysis confirms that the curated African NP seed set covers \qty{94.9}{\percent} of the African Natural Products Database (ANPDB, \num{9309} molecules) at Tanimoto similarity $\geq \num{0.4}$, indicating high coverage of major African NP collections. However, \qty{37.0}{\percent} of seed scaffolds (\num{91}/\num{246} Bemis--Murcko frameworks) are unique to our curated library and absent from ANPDB, showing that our hybrid NP + SD seed set contributes genuinely new chemical space beyond prior compilations.


\subsection*{African natural product chemical space as a source of antimalarial scaffolds}

The \num{13} NP-related optimization candidates (\qty{17.1}{\percent} of the \num{76} optimization-worthy centroids) validate African natural products as a productive source of antimalarial scaffolds, spanning \num{2} secondary hits (MPO $\geq \num{0.50}$) to \num{11} promising scaffolds (MPO \numrange{0.40}{0.50}). Their structural connection to African medicinal plants links computational discovery to traditional ethnopharmacological knowledge, and their identification through centroid-based sampling (\qty{99.3}{\percent} cost reduction) makes systematic exploration of African NP-derived chemical space accessible to malaria-endemic research groups. The \qty{17.1}{\percent} NP-relatedness rate is a deliberate consequence of exploring beyond the immediate seed vicinity — the \qty{82.9}{\percent} of optimization candidates representing genuinely new chemical space positions the framework to both exploit known bioactive scaffolds and explore adjacent chemotypes for new antimalarial candidates.

\subsection*{Named lead analysis}
\label{sec:discussion-leads}

Ligand 201 is the highest-affinity DHFR binder identified (Vina \qty{-8.49}{\kilo\calorie\per\mole}), with predicted interactions at Asp54 and Arg57. Ligand 438 is the sole \texttt{GOOD} hit against \emph{Pf}ATP4 (Vina \qty{-5.86}{\kilo\calorie\per\mole}). Polypharmacological candidates (164, 440, 169, 161) bind $\geq 2$ targets. The top-10 compounds share a piperazine scaffold (SA $< \num{2.0}$) with reductive amination as the predominant synthetic route. Detailed interaction diagrams and retrosynthetic analysis are presented in \Cref{SM-fig:interaction_leads,SM-tab:sm_s20_retrosynthesis}.

\subsection*{Limitations and future directions}

This computational study provides prioritized starting points for experimental validation but has inherent limitations. While \textit{in vitro} experimental assays were beyond the scope of this initial framework development, we established rigorous orthogonal computational-experimental proxies to justify candidate prioritization. The high early enrichment (EF\qty{5}{\percent}) in retrodicting known MMV Malaria Box ground-truth hits (overall hit rate \qty{69.8}{\percent}, $n = \num{399}$ confirmed actives) supports our screening protocol, and the synthetic accessibility threshold (SYBA $> \num{0}$) restricts our generative output to feasible molecules. Our evaluation strategy goes beyond the recent paradigm shift in generative chemistry benchmarking, exemplified by MolGenBench \cite{Cao2025molgenbench} and Durian \cite{Nie2024durian}, which departs from generic 1D property validities (GuacaMol \cite{Brown2019}) toward target-aware evaluation: our MMV Malaria Box retrodiction directly validates predicted binding against a library of \num{399} compounds with experimentally confirmed antimalarial activity, establishing a biologically grounded benchmark that precedes these community initiatives. Predicted retrosynthetic routes for the top-10 candidates (\Cref{SM-tab:sm_s20_retrosynthesis}) were generated via the ASKCOS public platform (MIT) and identify reductive amination as the primary synthetic strategy; however, these routes are computational predictions based on the Reaxys template library and require experimental verification. We have archived all \num{19913} synthesizable primary leads at Zenodo (DOI: 10.5281/zenodo.19608875) to facilitate experimental follow-up by malaria research programs with synthesis capacity. Nevertheless, all activity and ADMET predictions are computational estimates derived from machine learning models, and their prospective accuracy for African natural products requires physical experimental verification.

External calibration against the Tartarus docking benchmark \citep{nigam2023tartarus} (QuickVina/Smina on three standardized targets) confirms the independence of our MPO scoring dimensions. An initial 50-molecule sample yielded no correlation between the Tartarus composite binding affinity and the weighted MPO score (Spearman $\rho = \num{-0.037}$, $p = \num{0.777}$). A subsequent full-scale validation on \num{19913} molecules (\num{17211} with valid composite scores after Lipinski/substructure filtering) replicated this result with \num{280}$\times$ greater statistical power: Tartarus composite vs.\ weighted MPO, Spearman $\rho = \num{0.013}$, $p = \num{0.091}$ (non-significant). Neither the full-run composite nor any per-target Tartarus score ($|\rho| < \num{0.15}$ across all three targets) correlates meaningfully with our MPO ranking, confirming that drug-likeness and binding affinity are orthogonal properties in this chemical space. This orthogonality is a deliberate design feature motivating the combined MPO + docking ranking adopted in our candidate selection: ranking by MPO alone would prioritize drug-like molecules that may not bind, while ranking by docking alone would prioritize tight binders that may not be developable. The combined approach ensures that top-ranked compounds satisfy both drug-likeness and predicted potency criteria, a balance that neither single-criterion screening can achieve. The docking protocol employs rigid receptor approximations, consistent with standard virtual screening practices \citep{kitchen2004, shoichet2004}, and may not capture induced-fit effects or protein flexibility; a companion molecular dynamics study (Paper~2) validates top candidates against six clinically relevant resistance mutations (PfDHFR N51I/C59R/S108N/I164L; PfCRT K76T/K76A) via 300~ns total simulations to address this limitation. The MPO framework weights reflect general drug discovery priorities rather than malaria-specific optimization, and future work could refine these weights using experimental IC$_{50}$ data from lead compounds.

The dataset expansion strategy (Cheese API + STONED-SELFIES) was designed to balance exploration and scaffold preservation, but parameter choices (Tanimoto threshold \num{0.7}, \num{100} variants per seed) represent pragmatic selections rather than systematically optimized values. The low Tanimoto similarity (mean \num{0.206}) combined with high scaffold recovery (\qty{69.3}{\percent}) indicates successful substituent exploration while maintaining core pharmacophores, but experimental validation is needed to confirm that bioactivity-relevant features were preserved during expansion.

Future work should prioritize experimental validation of top-ranked compounds through \textit{in vitro} assays against \textit{Pf} 3D7 and drug-resistant strains, followed by cytotoxicity assessment and selectivity profiling. Synthesis feasibility for the \num{19913} primary leads should be evaluated, and a subset of structurally diverse compounds should be selected for experimental testing to validate the computational predictions and refine the MPO framework.

Several additional limitations should be noted. First, while docking grid boxes were uniformly set to \qty{25}{\angstrom}$\times$\qty{25}{\angstrom}$\times$\qty{25}{\angstrom} for all targets, larger boxes (e.g., \qty{30}{\angstrom} cube) may be warranted to capture peripheral binding interactions for membrane transporters such as PfCRT and PfATP4. Second, protein protonation states were assigned at pH 7.4 without explicit pKa calculations for active-site residues, which may affect hydrogen-bond predictions; PfCRT functions in the acidic digestive vacuole (pH~5.0--5.4), and its protonation state is a specific concern. To quantify this impact, we re-protonated the PfCRT structure at pH 5.2 using pKa predictions from PROPKA3~\cite{Olsson2011} followed by PDBFixer~\cite{Eastman2017} and re-docked the top~100 PfCRT-targeted candidates. Eight active-site residues exhibited shifted protonation states at vacuolar pH (ASP271, pKa~6.27; ASP274, pKa~7.05; GLU49, pKa~7.35; GLU143, pKa~6.53; GLU177, pKa~6.10; HIS68, pKa~6.18; HIS125, pKa~5.78; HIS218, pKa~6.74). The ranking correlation between pH~7.4 and pH~5.2 docking scores was moderate (Spearman $\rho = \num{0.270}$, $p = \num{0.007}$), with a systematic shift of \qty{+2.2}{\kilo\calorie\per\mole} in mean affinity, consistent with fewer charge-assisted interactions under acidic conditions. This confirms that pH correction measurably alters PfCRT rankings and is recommended for future PfCRT-targeted campaigns. Other targets operate at neutral pH and are unaffected. Third, all docking was performed against single static crystal structures without accounting for protein conformational ensembles or resistance mutations (e.g., PfCRT K76T/K76A, PfDHFR N51I/C59R/S108N/I164L). A companion study (Paper~2) validates top candidates against these mutations via molecular dynamics simulations, introducing a Resistance Resilience Score (RRS) for prioritization. Fourth, centroid-based sampling may not fully capture activity cliffs within clusters where minor structural changes produce large potency differences; the twofold intra- over inter-cluster Tanimoto separation (mean \num{0.227} vs \num{0.101}) indicates coherent clustering, but we recommend full-cluster rescoring for top-priority clusters during hit-to-lead optimization. We performed this rescoring for the top-20 MPO clusters (\num{1815} molecules) and confirmed smooth landscapes ($\sigma \leq \qty{0.38}{\kilo\calorie\per\mole}$) with no activity cliffs detected, while the best member consistently outperformed its centroid ($\Delta$ up to \qty{-0.66}{\kilo\calorie\per\mole}), confirming the approach's validity. Fifth, the \num{19913} synthesizable leads require experimental IC$_{50}$ testing before any biological conclusions can be drawn. Sixth, polypharmacology predictions ($>\num{70}$ multi-target candidates) are computational hypotheses that require orthogonal experimental confirmation (enzyme assays, cellular phenotypic screens). Seventh, the prospective DEKOIS external validation (PfDHFR, Vina only) achieved ROC-AUC \num{0.509}---near-random enrichment---confirming that AutoDock Vina docking alone cannot discriminate actives from property-matched decoys \cite{Bauer2013}. The consensus MMV enrichment (ROC-AUC \numrange{0.924}{1.000}) incorporates DiffDock ML rescoring, which provides the critical enrichment gain; however, the MMV benchmark remains a positive-control (in-distribution) rather than a fully independent decoy set.Prospective experimental validation of top-ranked compounds against \textit{Pf} 3D7 and drug-resistant strains is therefore essential to confirm the protocol's discriminative power.

Eighth, the PfCRT structure used for docking (PDB: 6UKJ) corresponds to the 3D7 isoform; the chloroquine-resistant 7G8 isoform prevalent in field isolates differs by four key mutations (C72S, K76T, A220S, N326D) that significantly alter the drug-binding pocket geometry. Docking against the 3D7 isoform may not fully recapitulate the binding landscape of the 7G8 variant, and prospective studies should consider isoform-specific docking. Similarly, the PfATP4 docking was performed against an apo-like inward-facing conformation (PDB: 9N10), which may not capture the outward-facing (E2) state critical for ATP-dependent ion efflux. The E1/E2 conformational cycle presents substantially different binding cavities at the membrane-cytosol interface, and our static apo-structure docking cannot account for this conformational plasticity. A companion molecular dynamics study (Paper~2) partially addresses these limitations through membrane-embedded simulations of the top candidates, capturing induced-fit stabilization, but prospective experimental validation of the predicted binding modes remains essential.

\section{Conclusion}

This study establishes a resource-efficient computational framework for antimalarial drug discovery by integrating deep generative chemistry with a dual-filter consensus docking and multi-parameter optimization (MPO) approach. By employing variational autoencoder (VAE) clustering for centroid-based sampling and Metropolis-Hastings MCMC for latent space optimisation, we substantially reduced virtual screening costs while maintaining high predictive accuracy on the MMV Malaria Box positive-control benchmark (consensus-score ROC-AUC \numrange{0.924}{1.000}) and demonstrated that local exploration of the VAE latent space yields measurable MPO improvement (\num{+0.0246}). Prospective DEKOIS PfDHFR docking (Vina only, ROC-AUC \num{0.509}) confirmed that structure-based docking alone cannot discriminate actives from property-matched decoys, validating the essential role of the ML-based DiffDock rescoring in achieving enrichment. This methodology enables the systematic exploration of massive molecular libraries (e.g., \num{65856} molecules) using standard hardware, addressing a fundamental computational bottleneck in the characterization of diverse chemical spaces.

Our results demonstrate the effectiveness of combining African natural product motifs with synthetic drug optimization through hybrid pharmacophore design. The identification of \num{28679} synthesizable leads (from \num{53} unique clusters, mean cluster size \num{763.8}) with favorable predicted selectivity and drug-likeness confirms that generative expansion can preserve privileged bioactive scaffolds while substantially diversifying chemical space. The DEKOIS external validation (AUC \num{0.509} for Vina only) versus the MMV consensus enrichment (AUC \numrange{0.924}{1.000}) demonstrates that the ML-based DiffDock rescoring provides the critical enrichment gain over naive docking, successfully identifying high-affinity candidates across essential \textit{P. falciparum} resistance targets including DHFR, CRT, and ATP4.

The accessibility of validated virtual screening protocols presented here offers a scalable solution for drug discovery in resource-limited settings. By reducing the reliance on massive high-performance computing infrastructure, this framework enables researchers in malaria-endemic regions to systematically prioritize local chemical space for antimalarial candidates. Ultimately, the open-source implementation and validated discovery of multiple polypharmacological leads provide a foundation for collaborative drug development efforts targeting drug-resistant malaria and other neglected tropical diseases.

% ACS-style back matter sections

\section*{Supporting information}
Supporting information available at the publisher's website. \Cref{SM-tab:scaffolds,SM-tab:selectivity,SM-tab:cyp450,SM-tab:stage_activity,SM-tab:pca,SM-tab:scaffold_recovery,SM-tab:fsp3,SM-tab:druglikeness,SM-tab:consensus_hits,SM-tab:source_breakdown,SM-tab:pca_variance,SM-tab:tanimoto_novelty,SM-tab:diffdock_vina_correlation,SM-tab:diversity_comparison,SM-tab:virtual_screening,SM-tab:sm_s16_admet_crossval,SM-tab:sm_s17_mpo_sensitivity,SM-tab:sm_s18_expansion_sensitivity,SM-tab:sm_s18b_chembl_enrichment,SM-tab:sm_s21_polypharm,SM-tab:sm_s20_retrosynthesis,SM-tab:mmv_validation,SM-tab:redocking_validation,SM-tab:consensus,SM-tab:computational_comparison,SM-tab:clustering_comparison,SM-tab:vae_comparison,SM-fig:vae_training,SM-fig:clustering_quality,SM-fig:enrichment_curves,SM-fig:interaction_leads,SM-fig:redocking_rmsd,SM-fig:consensus_comparison,SM-fig:consensus_correlation,SM-fig:sm_s15_top10_structures,SM-fig:si_umap_latent,SM-fig:si_phys_chem_clus}. MPO equations archived in repository.

\section*{Author information}

\subsection*{Corresponding author}
*E-mail: myke-vital.sao@facsciences-uy1.cm

\subsection*{Author contributions}
Conceptualization, M.V.S.T. and S.G.N.E.; methodology, M.V.S.T. and J.-P.T.N.; software, M.V.S.T.; validation, M.V.S.T., J.-P.T.N. and W.F.M.; formal analysis, M.V.S.T.; investigation, M.V.S.T.; resources, S.G.N.E. and W.F.M.; data curation, M.V.S.T.; writing--original draft preparation, M.V.S.T.; writing--review and editing, M.V.S.T., J.-P.T.N., P.S., W.F.M. and S.G.N.E.; visualization, M.V.S.T.; supervision, S.G.N.E. and W.F.M.; project administration, S.G.N.E. All authors read and agreed to the published version.

\subsection*{Funding}
This research received no external funding.

\subsection*{Notes}
The authors declare no competing financial interest.

\section*{Acknowledgments}
We acknowledge the Ersilia Open Source Initiative (\url{https://www.ersilia.io}, accessed on December 2025) for machine learning models, ChEMBL and DrugBank for compound data, MMV for Malaria Box, and RDKit.

\section*{Data availability}
SMILES strings for all \num{65856} compounds and primary lead structures are available in the Supplementary Materials. 3D ligand structures, computational workflows, and raw docking results have been deposited on Zenodo (DOI: 10.5281/zenodo.19608875) and are also archived in the project repository (\url{https://github.com/Vital-Sao/Malaria_codes}, accessed on 16 April 2026). The dataset is FAIR-compliant: \textbf{Findable} (permanent DOI via Zenodo), \textbf{Accessible} (freely downloadable without login), \textbf{Interoperable} (open formats: CSV, SMILES, PDB, PDBQT), and \textbf{Reusable} (CC BY 4.0 license). While raw enrichment metric CSVs are not archived as standalone files, the source datasets for docking threshold calibration (\texttt{p1\_threshold\_calibration.csv}) and MMV Malaria Box target consensus docking are fully provided in the repository to support reproducibility.

\bibliography{../bibliography/Sao_Chim_Space}

\end{document}

